% !TEX root = ../Main.tex
\chapter{条件概率}

"条件概率"回答：已知 $B$ 发生，$A$ 发生的概率是多少？它的操作是\textbf{限制样本空间}——把 $\Omega$ 缩至 $B$，然后把原事件 $A$ 替换为 $A \cap B$，重新计算比例。

\section{条件概率的定义}

\defn{条件概率}{
    若 $P(B) > 0$，事件 $A$ 在 $B$ 发生条件下的\textbf{条件概率}为
    \[
    P(A \mid B) = \frac{P(A \cap B)}{P(B)}.
    \]
    从古典概型视角：把样本空间限定到 $B$，计算 $A \cap B$ 与 $B$ 的比——"\textbf{缩小样本空间后再数}"。
}

\state{乘法公式}{
    由定义变形：
    \[
    P(A \cap B) = P(B) \cdot P(A \mid B) = P(A) \cdot P(B \mid A).
    \]
    对三事件：
    \[
    P(A \cap B \cap C) = P(A) \cdot P(B \mid A) \cdot P(C \mid A \cap B).
    \]
    即\textbf{按步骤分解}的乘法公式——"按步骤抽取"问题的自然表述。
}

\state{独立性的等价刻画}{
    当 $P(B) > 0$ 时，
    \[
    A, B \ \text{独立}\ \iff\ P(A \mid B) = P(A).
    \]
    直观："$B$ 发生的信息不改变 $A$ 的概率"——这正是"独立"的直觉说法。
}

\ex[乘法公式的步骤应用]{
    袋中有 $4$ 只红球、$3$ 只白球，不放回地依次取 $2$ 球。求第一次取到红球、第二次取到白球的概率。
}

\sol{
    设 $A$ = "第一次取红"，$B$ = "第二次取白"。

    $P(A) = \dfrac{4}{7}$。在 $A$ 发生条件下，袋中剩 $3$ 红 $3$ 白，共 $6$ 个，故 $P(B \mid A) = \dfrac{3}{6} = \dfrac{1}{2}$。

    \[
    P(A \cap B) = P(A) \cdot P(B \mid A) = \frac{4}{7} \cdot \frac{1}{2} = \frac{2}{7}.
    \]
}

\section{全概率公式}

\state{全概率公式}{
    若 $B_1, B_2, \ldots, B_n$ 是\textbf{样本空间的一个划分}（两两互斥、并起来是 $\Omega$，每个 $P(B_i) > 0$），则对任意事件 $A$：
    \[
    P(A) = \sum_{i=1}^n P(B_i) \cdot P(A \mid B_i).
    \]

    \textbf{直觉}：把 $A$ 按不同"原因" $B_i$ 拆开，各原因导致 $A$ 的贡献权重为 $P(B_i)$，条件贡献为 $P(A \mid B_i)$——加权求和即得 $P(A)$。
}

\ex[全概率]{
    某工厂有两台机器 $M_1, M_2$ 生产同一种零件，产量比 $6 : 4$。$M_1$ 的次品率为 $3\%$，$M_2$ 的次品率为 $5\%$。任取一件产品，求其为次品的概率。
}

\sol{
    设 $B_i$ = "来自 $M_i$"，$A$ = "次品"。$P(B_1) = 0.6$，$P(B_2) = 0.4$；$P(A \mid B_1) = 0.03$，$P(A \mid B_2) = 0.05$。
    \[
    P(A) = 0.6 \cdot 0.03 + 0.4 \cdot 0.05 = 0.018 + 0.020 = 0.038.
    \]
    即 $3.8\%$。
}

\rmkb{
    \textbf{条件概率的三种形式}：
    \begin{itemize}
        \item \textbf{定义式} $P(A | B) = \dfrac{P(AB)}{P(B)}$——由联合概率 $P(AB)$ 和边缘 $P(B)$ 求条件；
        \item \textbf{乘法式} $P(AB) = P(B) P(A | B)$——由条件概率和边缘反求联合；
        \item \textbf{全概率} $P(A) = \sum P(B_i) P(A | B_i)$——由"各原因"的贡献合成边缘。
    \end{itemize}

    三种形式\textbf{是同一公式的三种视角}：条件、联合、边缘概率之间的转换。哪个已知就用哪个——\textbf{"已知的当条件、未知的当目标"}是灵活使用的关键。
}
